\paragraph{Identifying the SFibAI baseline.}
The method source is \citet{xu2026sfibai}. The baseline follows the authors' public implementation. The baseline uses ResNet-50, global average pooling, a linear $2048\rightarrow36$ head, softmax, and an expected score on 0.0--3.5. View and localization prediction heads are absent. Upstream permits an explicit initialization checkpoint and otherwise constructs a randomly initialized backbone; our common ImageNet initialization is an adaptation for this comparison, not an upstream default.

\paragraph{Grading objectives.}
VCRE-Fib and SFibAI use the same fibrosis-score scale and hybrid-loss weights $(\alpha,\beta,\gamma)=(1,0.02,0.02)$. For target bin $j\in\{0,\ldots,35\}$ and predicted probabilities $p_k$, both convert bin indices to real-valued clinical scores before computing MSE and the clinical-grade boundary penalty:
\[
\begin{aligned}
y&=0.1j,\qquad \hat y=0.1\sum_{k=0}^{35}kp_k,\\
L_{\mathrm{grade}}&=L_{\mathrm{KL,local}}+0.02(\hat y-y)^2
 +0.02|\hat y-y|\mathbf{1}(\hat c\ne c).
\end{aligned}
\]
Here $c$ and $\hat c$ are obtained from $y$ and $\hat y$ using thresholds 0.5, 1.5, and 2.5, with a score exactly at a threshold assigned to the higher grade. The displayed per-image terms are averaged over the batch. Both MSE and the absolute-error boundary term are evaluated on the 0.0--3.5 score scale.

The local-KL term uses a normalized Gaussian soft target over 36 bins with standard deviation one bin and a local window centered on the target bin. Both methods use this hybrid grading-loss formulation with aligned score units and weights.

\paragraph{Training protocol.}
All four models use the same random seed (2026), ImageNet-1K-V2 weights, batch size 24, bfloat16 precision, AdamW learning rate $10^{-4}$ and weight decay $10^{-4}$, and StepLR factor 0.6 every 15 epochs. They share ROI preprocessing, stretch resize, right-angle rotation/saturation augmentation, patient-disjoint splits, and the grading objective above. This standardizes the public SFibAI release's different crop/augmentation and initialization defaults. Training trajectories and checkpoint comparison use the same 90-epoch cap for SFibAI, both deletion variants, and full VCRE-Fib. Validation selects from epochs 21--90 by $\risk$, then image COR and earlier epoch, replacing the upstream tolerance-accuracy selector. The selected epochs are 65, 80, 39, and 47, respectively. This comparison on the curated cohort is not a reproduction of the original SFibAI cohort experiment.

\paragraph{Additional VCRE-Fib supervision.}
Six-class view cross-entropy has weight 0.1. The weak-box loss combines inside log-sum-exp pooling (temperature 10) and the outside response fraction with unit inner weights and overall weight 0.1. Only valid nonempty boxes with target bin $j>0$ participate; exact score 0.0 is excluded, not the entire F0 grade. Weak targets are transformed in the ROI frame. The support map is learned from grading. Auxiliary losses update the shared encoder, while $q$ and $a$ are detached only at their grading interfaces. The w/o View variant removes the view head, six adapters, posterior mixture, and view CE while retaining unconditioned regional evidence and weak localization. The w/o Weak Loc.\ variant removes the localization head and box loss, using neutral $a_u=1$ with the view and support paths retained. Each deletion variant starts independently from the prescribed initialization.

